\section{Условия выполнения программы}

\subsection{Минимальный состав технических средств}

Для выполнения программы требуется вычислительное устройство, обеспечивающее
запуск 64-разрядных приложений и запись выходных файлов на локальную файловую
систему либо в доступное оператору хранилище.

Минимально необходимыми техническими средствами являются:
\begin{enumerate}
    \item 64-разрядный процессор;
    \item оперативная память, достаточная для формирования выходных батчей;
    \item дисковое пространство, достаточное для размещения результатов
    генерации.
\end{enumerate}

\subsection{Максимальный состав технических средств}

При эксплуатации программы на больших значениях масштаба генерации
целесообразно использовать многоядерный процессор, увеличенный объем
оперативной памяти и дисковую подсистему с высокой пропускной способностью.
Это позволяет эффективно использовать многопоточную обработку партиций и
сократить время записи выходных данных.

\subsection{Программные средства}

Для выполнения программы должны быть доступны:
\begin{enumerate}
    \item 64-разрядная операционная система семейства Linux;
    \item собранный исполняемый файл программы \texttt{schengen\_main};
    \item программные библиотеки, необходимые для форматов, включенных в
    текущую сборку программы;
    \item права на чтение входного каталога запуска и запись в каталог или URI
    вывода.
\end{enumerate}

\subsection{Дополнительные условия выполнения программы}

При использовании удаленных URI оператор должен обеспечить сетевую доступность
соответствующей файловой системы или объектного хранилища.

При выборе форматов вывода оператор должен учитывать, что набор поддерживаемых
форматов определяется конфигурацией конкретной сборки программы.
